\section{The complete law of an arbitrary symplectic compression}
\label{sec:complete-law}

Let $q$ be a prime power. All spaces and matrices in this section are over
$\mathbb F_q$. An alternating form means a bilinear form $B$ satisfying
$B(x,x)=0$ for every $x$. Its matrix therefore has zero diagonal as well as
$H^{\mathsf T}=-H$; both requirements matter in characteristic two.
Let $V$ be a $2n$-dimensional space with a nondegenerate alternating form
$\langle\ ,\ \rangle$. Write $J$ for its matrix in a chosen basis.

Fix two injective linear maps
\[
 Y:\mathbb F_q^a\longrightarrow V,\qquad
 X_0:\mathbb F_q^b\longrightarrow V,
\]
and let
\[
 H_1=Y^{\mathsf T}JY,\quad H_0=X_0^{\mathsf T}JX_0,
 \qquad \operatorname{rank}H_i=2s_i.
\]
For $g$ chosen uniformly from $\operatorname{Sp}(V)$, the observed matrix is
\begin{equation}\label{eq:complete-observation}
 M=Y^{\mathsf T}JgX_0\in\operatorname{Mat}_{a,b}(\mathbb F_q).
\end{equation}
Neither prescribed form need be zero or nondegenerate. In particular, both
observed subspaces may have arbitrary nonzero radicals. The necessary and
sufficient conditions that the two initial frames exist are
$a-s_1\le n$ and $b-s_0\le n$.

\subsection{An explicit count for a prescribed Gram matrix}

For integers $d,s\ge0$, define
\begin{equation}\label{eq:E-count}
 \mathcal E_n(d,s)=
 q^{2nd-\binom d2}
 \prod_{j=n-d+s+1}^{n}(1-q^{-2j})
 \quad\text{if }0\le2s\le d\text{ and }d-s\le n,
\end{equation}
and set $\mathcal E_n(d,s)=0$ otherwise. Empty products are one; in
particular $\mathcal E_n(0,0)=1$. Formula \eqref{eq:E-count} is only
evaluated in the stated feasible range, so it contains no nonpositive
product indices there.

\begin{lemma}[Prescribed-Gram frames]\label{lem:E-count}
Let $E$ be a specified $d$-dimensional alternating space of rank $2s$.
The number of isometric injections $E\longrightarrow V$ is
$\mathcal E_n(d,s)$. The symplectic group acts transitively on this set
whenever it is nonempty.
\end{lemma}

\begin{proof}
Decompose $E=S\oplus R$, where $S$ is nondegenerate of dimension $2s$
and $R$ is the radical, of dimension $h=d-2s$. A symplectic basis of $S$
embeds in $V$ in
\[
 \frac{|\operatorname{Sp}_{2n}(q)|}
 {|\operatorname{Sp}_{2(n-s)}(q)|}
\]
ways. Once its image is fixed, the images of an ordered basis of $R$
must be an independent isotropic $h$-frame in the orthogonal complement,
which is symplectic of dimension $2(n-s)$. After $i$ radical vectors have
been chosen, the next one lies in their orthogonal complement and outside
their span. There are exactly
$q^{2(n-s)-i}-q^i$ choices. Hence the count is
\begin{equation}\label{eq:E-integer}
 \frac{|\operatorname{Sp}_{2n}(q)|}
 {|\operatorname{Sp}_{2(n-s)}(q)|}
 \prod_{i=0}^{d-2s-1}(q^{2(n-s)-i}-q^i),
\end{equation}
and it is nonzero exactly when $h\le n-s$.

The standard symplectic-basis count gives
$|\operatorname{Sp}_{2u}(q)|=q^{u^2}\prod_{j=1}^u(q^{2j}-1)$.
Substituting in \eqref{eq:E-integer} and extracting powers of $q$ gives
\eqref{eq:E-count}; the remaining factors have consecutive indices
$n-d+s+1,\ldots,n$.

For transitivity, first extend the image of the chosen symplectic basis
of $S$ to a symplectic decomposition of $V$. In its complement, extend
the ordered isotropic basis of $R$ to a symplectic basis by adjoining
dual vectors and then a basis of the residual complement. The change of
symplectic basis between any two such completions sends one entire
prescribed frame to the other. This is the relevant elementary case of
Witt extension, and the argument works in characteristic two.
\end{proof}

\subsection{Three invariants of an observed matrix}

For a proposed value $A\in\operatorname{Mat}_{a,b}(\mathbb F_q)$, form
the alternating matrix
\begin{equation}\label{eq:combined-Gram}
 \Gamma_A=
 \begin{pmatrix}H_1&A\\-A^{\mathsf T}&H_0\end{pmatrix}.
\end{equation}
Write $\operatorname{rank}\Gamma_A=2s$ and
$r=a+b-2s$. In $E=E_1\oplus E_0=\mathbb F_q^a\oplus\mathbb F_q^b$,
let $R=\operatorname{rad}\Gamma_A$ and put
\begin{align}
 t_1&=\dim(R\cap E_1)=\dim(\ker H_1\cap\ker A^{\mathsf T}),
 \label{eq:t1}\\
 t_0&=\dim(R\cap E_0)=\dim(\ker H_0\cap\ker A),
 \label{eq:t0}\\
 c&=r-t_0-t_1.\label{eq:complement-c}
\end{align}
Thus $r$ and $c$ are determined by the three integers $(s,t_0,t_1)$
and the fixed frame dimensions.

Write $\genfrac{[}{]}{0pt}{}{u}{v}_q$ for the number of $v$-subspaces
of $\mathbb F_q^u$, with value zero outside $0\le v\le u$, and set
$|\operatorname{GL}_0(q)|=1$. For $0\le\ell\le r$ define
\begin{equation}\label{eq:T-count}
 \begin{split}
 T_\ell(r,t_0,t_1)=
 \sum_{k=\max(0,\ell-c)}^{\min(\ell,t_0,t_1)}
 &\genfrac{[}{]}{0pt}{}{t_0}{k}_q
 \genfrac{[}{]}{0pt}{}{t_1}{k}_q
 |\operatorname{GL}_k(q)|
 \genfrac{[}{]}{0pt}{}{c}{\ell-k}_q\\[-2pt]
 &\hspace{12mm}\cdot q^{(\ell-k)(t_0+t_1-k)}.
 \end{split}
\end{equation}
An empty sum is zero. This is a nonnegative integer formula with no
cancellation.

\begin{lemma}[Relations meeting neither frame]\label{lem:transverse-relations}
The number of $\ell$-dimensional subspaces $L\le R$ satisfying
$L\cap E_0=L\cap E_1=0$ is $T_\ell(r,t_0,t_1)$.
Such a subspace exists exactly when
\begin{equation}\label{eq:ell-support}
 0\le\ell\le r-\max(t_0,t_1)=c+\min(t_0,t_1).
\end{equation}
\end{lemma}

\begin{proof}
Put $R_i=R\cap E_i$, choose a complement $C$ to $R_0\oplus R_1$
in $R$, and set $J_L=L\cap(R_0\oplus R_1)$. Both coordinate
projections of $J_L$ are injective. If $\dim J_L=k$, it is therefore
the graph of an isomorphism between a $k$-subspace of $R_0$ and a
$k$-subspace of $R_1$. The number of choices is
\[
 \genfrac{[}{]}{0pt}{}{t_0}{k}_q
 \genfrac{[}{]}{0pt}{}{t_1}{k}_q|\operatorname{GL}_k(q)|.
\]
The image $D$ of $L$ under projection to $C$ has dimension $\ell-k$,
giving $\genfrac{[}{]}{0pt}{}{c}{\ell-k}_q$ choices.
For fixed $J_L,D$, the possible quotients $L/J_L$ are exactly the
graphs of linear maps
$D\longrightarrow(R_0\oplus R_1)/J_L$, of which there are
$q^{(\ell-k)(t_0+t_1-k)}$. They all have the required intersection
with $R_0\oplus R_1$ and meet neither $R_i$. This proves
\eqref{eq:T-count}.

The dimension inequality $\ell+\dim R_i\le r$ is necessary.
It is sufficient: take the graph of an isomorphism between subspaces
of $R_0,R_1$ of dimension $\min(t_0,t_1)$, add $C$, and then take
any desired lower-dimensional subspace. This proves \eqref{eq:ell-support}.
\end{proof}

\begin{theorem}[Complete compression law]\label{thm:complete-law}
For the observation \eqref{eq:complete-observation} and every proposed
matrix $A$, its probability is
\begin{equation}\label{eq:main-law}
 \boxed{\displaystyle
 \Pr(M=A)=
 \frac{\displaystyle\sum_{\ell=0}^{r}
 T_\ell(r,t_0,t_1)\,\mathcal E_n(a+b-\ell,s)}
 {\mathcal E_n(a,s_1)\,\mathcal E_n(b,s_0)}.}
\end{equation}
All quantities on the right are determined by $(s,t_0,t_1)$ and the
fixed parameters. This is valid over every finite field.
More precisely, the $\ell$-th summand divided by the same denominator
is
\begin{equation}\label{eq:joint-law}
 \Pr\bigl(M=A,\ \dim(\operatorname{im}Y\cap
 \operatorname{im}(gX_0))=\ell\bigr).
\end{equation}
\end{theorem}

\begin{proof}
By Lemma~\ref{lem:E-count}, $gX_0$ is uniform among the
$\mathcal E_n(b,s_0)$ prescribed-Gram injections. The number of
such $X$ giving a specified $A$ is independent of the particular
target injection $Y$: a symplectic change of coordinates carrying
one target to another carries the corresponding source injections
bijectively. The same holds if the intersection dimension is also
specified. We may consequently choose $Y$ independently and uniformly
among its $\mathcal E_n(a,s_1)$ possible values and count pairs.

Each pair $(Y,X)$ with $Y^{\mathsf T}JX=A$ is the linear map
\[
 F:E_1\oplus E_0\longrightarrow V,\qquad F(u,v)=Yu+Xv,
\]
with Gram matrix $\Gamma_A$, whose restrictions to $E_1,E_0$ are
injective. Its kernel $L$ lies in $R$, since
$\Gamma_A(z,w)=\langle Fz,Fw\rangle$ vanishes for every $w$ when
$Fz=0$. Injectivity on the two summands says exactly
$L\cap E_1=L\cap E_0=0$.

Conversely, fix any such $L$ of dimension $\ell$. The alternating
form descends to $E/L$ with dimension $a+b-\ell$ and rank $2s$.
Every isometric injection $E/L\longrightarrow V$ gives exactly one
map $F$ of the required kind with kernel exactly $L$, and hence one
pair $(Y,X)$. Lemma~\ref{lem:E-count} counts these injections by
$\mathcal E_n(a+b-\ell,s)$; Lemma~\ref{lem:transverse-relations}
counts possible kernels by $T_\ell$. Summing and dividing by the
number of independent pairs gives \eqref{eq:main-law}.

Finally, the map $(u,v)\mapsto Yu=-Xv$ identifies $\ker F$ with
$\operatorname{im}Y\cap\operatorname{im}X$. Thus the same count
at a fixed $\ell$ proves \eqref{eq:joint-law}.
\end{proof}

\subsection{The exact obstruction and the smallest possible ambient space}

\begin{theorem}[Exact support and minimal dimension]\label{thm:minimum-dimension}
In Theorem~\ref{thm:complete-law},
\begin{equation}\label{eq:support-criterion}
 \boxed{\quad\Pr(M=A)>0\quad\Longleftrightarrow\quad
 s+\max(t_0,t_1)\le n.\quad}
\end{equation}
More generally, start with any two alternating forms $H_1,H_0$ and
any cross-pairing matrix $A$, without assuming that frames have already
been realized. The smallest dimension of a nondegenerate alternating
space admitting injective realizations of both spaces, preserving
their forms and the prescribed cross pairings, is
\begin{equation}\label{eq:minimum-dimension}
 \boxed{\qquad 2\bigl(s+\max(t_0,t_1)\bigr).\qquad}
\end{equation}
\end{theorem}

\begin{proof}
A summand in \eqref{eq:main-law} is positive precisely when a
transverse kernel of dimension $\ell$ exists and its quotient embeds.
By the preceding lemmas, these conditions are
\[
 0\le\ell\le r-\max(t_0,t_1),\qquad
 \ell\ge a+b-s-n.
\]
Their compatibility is exactly $s+\max(t_0,t_1)\le n$, because
$r=a+b-2s$. The largest allowed kernel has dimension
$r-\max(t_0,t_1)$; its quotient has rank $2s$ and radical dimension
$\max(t_0,t_1)$, and embeds in a symplectic space of dimension
$2(s+\max(t_0,t_1))$ by Lemma~\ref{lem:E-count}.
The quotient map is injective on each original summand. Conversely,
any realization produces a transverse kernel and therefore cannot
use a smaller ambient dimension. This argument also proves
\eqref{eq:minimum-dimension} without separately assuming that the
two original frames are realizable.
\end{proof}

\subsection{Sufficient statistics and special cases}

The triple $(s,t_0,t_1)$ is sufficient for every likelihood test of
the symplectic law against uniform independent entries: both laws
are uniform within each set of matrices sharing a triple. Consequently
their full total variation distance is exactly the distance between
the induced distributions of that triple. To see this directly, let
$\mathcal C_\tau$ be a triple class and $p_\tau$ the value of
\eqref{eq:main-law} on it. Both distances are
\[
 \frac12\sum_\tau |\mathcal C_\tau|
 |p_\tau-q^{-ab}|.
\]
This does not assert that the triple is the coarsest possible statistic;
different triples can happen to give the same likelihood.

Conditional on an observed $A$ with positive probability, the
intersection dimension has the distribution
\begin{equation}\label{eq:conditional-intersection}
 \Pr(\ell\mid A)=
 \frac{T_\ell(r,t_0,t_1)\mathcal E_n(a+b-\ell,s)}
 {\sum_jT_j(r,t_0,t_1)\mathcal E_n(a+b-j,s)}.
\end{equation}
Conditional also on $\ell$, all allowed relation subspaces
$L\le R$ are equally likely, since the number of embeddings of
their quotients is the same.

If both $H_i$ are nondegenerate, then $t_0=t_1=0$ and
$T_\ell=\genfrac{[}{]}{0pt}{}{r}{\ell}_q$. The likelihood depends
only on the rank of $\Gamma_A$. For nondegenerate $H_1$, ordinary
block elimination gives the useful identity
\begin{equation}\label{eq:Schur-rank}
 \operatorname{rank}\Gamma_A
 =a+\operatorname{rank}(H_0+A^{\mathsf T}H_1^{-1}A).
\end{equation}
Thus the relevant rank is that of an alternating residual form,
which is generally different from the ordinary matrix rank of $A$.

At the previously studied Lagrangian-output endpoint, $a=n$ and
$H_1=0$. Put $K=\ker A$. A basis adapted to $A$ yields
\[
 \operatorname{rank}\Gamma_A
 =2\operatorname{rank}A+\operatorname{rank}(H_0|_{K\times K}),
 \qquad t_1=n-\operatorname{rank}A.
\]
Since $t_0\le\dim K=b-\operatorname{rank}A$, the general support
criterion directly gives the familiar isotropic-kernel obstruction
when also $b=n$: the proposed matrix occurs exactly when
$H_0|_{K\times K}=0$. The full theorem permits arbitrary dimensions
and arbitrary forms on both frames, so it reaches beyond that
endpoint.

The computational evaluation uses Gaussian elimination for
\eqref{eq:combined-Gram}--\eqref{eq:t0}, followed by at most
$O((a+b)^2)$ terms of scalar integer arithmetic. This is an operation
count, not a bit-complexity bound; the exact integers can be large.
